%% title: The test fold is inside a foundation model's training set
%% target: Sections/Results.tex
%% anchor: Open-set results are therefore reported with the rule named, and the rule is re-selected whenever the model or the domain changes.
%% action: after
%% why: Two thirds of the test fold is in BioCLIP-2's training data (by GBIF occurrence id, not species); a frozen probe understates it by 7 points; after adaptation neither it nor our best model dominates, so the contribution is the recipe, not a backbone.
%% source: paper/DRAFT.md §4.14; journal/2026-10-22-rescoring-in-the-mini-metrics-convention.md (P3c minus baseline +1.25 old, +1.32 zero, +1.34 calibrated; decontaminated fold 217,856 rows); data/rescore/probe-grouped.csv, probe-ho-grouped.csv (base 0.2875/0.2394, P3c 0.3768/0.3017, P1b 0.4108/0.2138, P4 0.4796/0.3568, P5 0.6223/0.4974, P5b 0.6181/0.4719, B8 0.6472/0.4520)
%% check: Paste after step 17. | The decontaminated fold is 219,048 images in DRAFT §4.14.1 and 217,856 rows in the rescoring; find which is right before pasting. | Learning rates and the frozen-probe numbers are in the old convention (identical on the decontaminated fold to within 1.3 points). | Asger's ViT-L/16 backbone: check its pretraining data for the same overlap.
%% status: ready
\subsection{The test fold is inside a foundation model's training set}\label{sec:results_contamination}
Global Lepidoptera and TreeOfLife-200M, the training set of BioCLIP-2 \textbf{(REF)}, are both drawn from GBIF.
Joined on the GBIF occurrence identifier, which identifies the photograph rather than the species, 93.3\,\% of our species, 65.4\,\% of our images and 413,865 of the 629,742 test images (65.7\,\%) are in BioCLIP-2's training data.
Every comparison with BioCLIP-2 on the test fold is therefore re-scored on the 35\,\% of the fold without an occurrence identifier in TreeOfLife-200M (\textbf{219,048} images), for both arms.
A taxonomic check, whether the model has seen these species, would have reported the same overlap as harmless; the identifier join reads only provenance columns and runs on a laptop.
The trap images carry no GBIF identifiers and are not contaminated.

On the clean part of the fold, a frozen cosine probe on BioCLIP-2 scores 0.844, 5.8 points below our baseline (0.902).
Fine-tuned, it scores 0.891, 0.903 and 0.915 at learning rates $10^{-3}$, $10^{-4}$ and $10^{-5}$: unfreezing is worth 7 points, and recipe L's default learning rate is the worst of the three.
On the trap images the fine-tuned model leads our baseline by 8.9 points on the probe set (0.377 against 0.288) and 6.2 on its held-out species.
Adapted to the traps with its backbone frozen, it reaches 0.480; unfrozen (P5), 0.622 on the probe set and 0.497 on held-out species, against 0.647 and 0.452 for B8, our best recipe-L model.
B8 leads by 2.5 points on the probe set and P5 by 4.5 on held-out species, where two identical P5 trainings differ by 2.5: neither dominates.
A representation trained on 43 times more images of the same domain is better before adaptation, and the adaptation recipe closes the difference.
